\documentclass[UTF8]{ctexart}
\usepackage{booktabs}
\usepackage{amsmath}
\usepackage{natbib}
\usepackage[margin=2.5cm]{geometry}

\title{中文短文的公式与引用检查}
\author{示例作者}
\date{}

\begin{document}
\maketitle

\begin{abstract}
本文用一段中文短文检验公式编号、交叉引用与文献引用的完整性。
\end{abstract}

\section{引言}
\label{sec:intro}
检索增强的问答系统把外部索引与小型阅读器结合
\citep{lewis2020rag}。稠密检索是该任务上常用的基线
\citep{karpukhin2020dpr}。中文场景下的评测口径见 \citet{zhao2019zh}。

\section{方法}
\label{sec:method}
给定问题 $q$ 与前 $k$ 个段落，阅读器对每个段落打分。训练目标为
\begin{equation}
  \label{eq:loss}
  \mathcal{L} = \frac{1}{N} \sum_{i=1}^{N} \ell(y_i, \hat{y}_i)
\end{equation}
其中 $N$ 是样本数。准确率定义为
\begin{equation}
  \label{eq:acc}
  \mathrm{Acc} = \frac{TP + TN}{TP + TN + FP + FN}
\end{equation}
式~\ref{eq:loss} 与式~\ref{eq:acc} 给出训练与评价的口径。

\section{结果}
\label{sec:results}
在合成开发集上，融合阅读器的准确率为 72.4 点，比稀疏基线高 11.2 点
\citep{lewis2020rag}。评价按式~\ref{eq:acc} 的定义计算。

\bibliographystyle{plainnat}
\bibliography{refs}

\end{document}
